%% This file contains a number of tweaks that are typically applied to the main document.
%% They are not enabled by default, but can be enabled by uncommenting the relevant lines.


%%
%% Official support for 2-column teaser figure
%%
\usepackage{cuted} %< for strip environment

%%
%% For auto-referencing of figures (see fig/teaser.tex and \label{} therein) 
%%
\usepackage{currfile} %< for \currfilebase macro

%%
%% caption and (global) spacing overrides
%%
\usepackage{caption} %< for "\captionof" commands (teaser)
% \setlength{\abovecaptionskip}{.5em} % space above caption
% \setlength{\belowcaptionskip}{1em} % space below caption

%%
%% Inline annotations; for predefined colors, refer to "dvipsnames" in the xcolor package:
%% https://tinyurl.com/overleaf-colors
%%
\newcommand{\red}[1]{{\color{red}#1}}
\newcommand{\todo}[1]{{\color{red}#1}}
\newcommand{\TODO}[1]{\textbf{\color{red}[TODO: #1]}}
%%
%% Disable for camera ready / submission by uncommenting these lines:
%%
% \renewcommand{\TODO}[1]{}
% \renewcommand{\todo}[1]{#1}

%%
%% Pointers into the supplement.  The supplement is a separate file (due 11-23), so the main text cannot
%% \cref into it.  \supp{label} prints "the supplement"; with \withsupptrue (flip it below, only for a draft
%% that compiles sec/X_suppl in) it also prints the section, so that every pointer can be checked.
%%
\newif\ifwithsupp
\withsuppfalse
\newcommand{\supp}[1]{the supplement\ifwithsupp~(\cref{#1})\fi}

%%
%% Work harder in optimizing text layout. Typically shrinks text by 1/6 of page, enable
%% it at the very end of the writing process, when you are just above the page limit.
%%
% \usepackage{microtype}

%%
%% Heading gaps, rigid (no plus/minus): sections 1.5 × the template default; paragraphs 3.11ex (1.5 × less 25%, less 15%; user 11-16).
%% 3.11ex is 0.96 × the article default's natural size and inside its own shrink range (3.25ex minus .2ex = 3.05ex): do NOT go below 3.05ex.
%%   \section   cvpr.sty: -10pt / 7pt     → -15pt / 10.5pt
%%   \paragraph article:  3.25ex plus 1ex minus .2ex → 3.11ex fixed
%% \parskip kept at 0pt so twocolumn \flushbottom cannot re-stretch gaps apart.
%% Run-in/@startsection kept; do not edit cvpr.sty (RED_FILE §1.3).
%%
\setlength{\parskip}{0pt}
\makeatletter
\def\cvprsection{\@startsection {section}{1}{\z@}%
  {-15pt}{10.5pt} {\large\bf}}
\renewcommand{\paragraph}{\@startsection{paragraph}{4}{\z@}%
  {3.11ex}%
  {-1em}%
  {\normalfont\normalsize\bfseries}}
\makeatother

%%
%% Globally adjusts space between figure and caption.
%%
% \setlength{\abovecaptionskip}{.5em}

%%
%% Allows "the use of \paper to refer to the project name"
%% with automatic management of space at the end of the word.
%%
% \usepackage{xspace}
% \newcommand{\paper}{ProjectName\xspace}

%%
%% Commonly used math definitions.
%%
% \DeclareMathOperator*{\argmin}{arg\,min}
% \DeclareMathOperator*{\argmax}{arg\,max}

%%
%% Tighten underline.
%%
% \usepackage{soul}
% \setuldepth{foobar}
